% GENERATED FILE. Edit canonical lessons, then run npm run generate:books.
% canonical-source-hash: fnv1a64:7a53ffcd6078cd45
% canonical-lessons: KA-C355-utkhanana, KA-C355-puratatva, KA-C355-avasesa, KA-C355-sasana, KA-C355-hastaprati

\chapter{Digging Up the Past}
\label{ch:ka-digging-up-the-past}
\hlchaptermodality{355}

\begin{chapteropening}
\textbf{By the end of this chapter:} \emph{I can say an excavation, archaeology, ruins, an inscription and a manuscript.}

Complete the last lesson of chapter 355: I can say an excavation, archaeology, ruins, an inscription and a manuscript.
\end{chapteropening}

\section[\texorpdfstring{\kn{ಉತ್ಖನನ} (utkhanana)}{ಉತ್ಖನನ (utkhanana)}]{\kn{ಉತ್ಖನನ} (utkhanana) — an excavation}

\label{lesson:KA-C355-utkhanana}



\begin{quote}
Before the new one: say the Kannada for a royal court, then the Kannada for a prince.
\end{quote}

\subsection*{You'll want to know: \kn{ಉತ್ಖನನ}}
\textbf{\kn{ಉತ್ಖನನ}} — \emph{utkhanana} — \textquotedblleft{}an excavation\textquotedblright{}.

\begin{quote}
\kn{ಉತ್ಖನನದಲ್ಲಿ} \kn{ಹಳೆಯ} \kn{ಮಡಕೆಗಳು} \kn{ಸಿಕ್ಕಿದವು}.
\end{quote}

\begin{quote}
\emph{utkhananadalli haḷeya maḍakegaḷu sikkidavu.}
\end{quote}

\begin{quote}
Old pots were found in the excavation.
\end{quote}

The careful digging done to find old things.

\begin{cousinweb}
\textbf{\kn{ಉತ್ಖನನ}} (\emph{utkhanana}) is borrowed from Sanskrit \emph{utkhanana}, \textquotedblleft{}digging or tearing out\textquotedblright{} (Monier-Williams).
\end{cousinweb}

\begin{grammarlens}[title={What you've built}]
One more word for telling what happened.
\end{grammarlens}

\subsection*{Your turn}
\begin{itemize}
\raggedright
  \item \emph{Say it:} \emph{utkhanana}
  \item \emph{Say it:} \emph{utkhanana}, once more
  \item \emph{Say it:} \emph{utkhananadalli haḷeya maḍakegaḷu sikkidavu.} — Old pots were found in the excavation.
  \item \emph{Recall:} say the Kannada for a royal court, then the Kannada for a prince, then say \emph{utkhanana} again
\end{itemize}

\begin{tcolorbox}[breakable,colback=teal!4,colframe=teal!35!black,title={Before you move on}]
What does \kn{ಉತ್ಖನನ} mean? (\textquotedblleft{}An excavation\textquotedblright{}.) Say it once more.
\end{tcolorbox}
\section[\texorpdfstring{\kn{ಪುರಾತತ್ವ} (purātatva)}{ಪುರಾತತ್ವ (purātatva)}]{\kn{ಪುರಾತತ್ವ} (purātatva) — archaeology}

\label{lesson:KA-C355-puratatva}



\begin{quote}
Before the new one: say the Kannada for a prince, then the Kannada for an excavation.
\end{quote}

\subsection*{You'll want to know: \kn{ಪುರಾತತ್ವ}}
\textbf{\kn{ಪುರಾತತ್ವ}} — \emph{purātatva} — \textquotedblleft{}archaeology\textquotedblright{}.

\begin{quote}
\kn{ಅವಳು} \kn{ಕಾಲೇಜಿನಲ್ಲಿ} \kn{ಪುರಾತತ್ವ} \kn{ಓದುತ್ತಿದ್ದಾಳೆ}.
\end{quote}

\begin{quote}
\emph{avaḷu kālējinalli purātatva ōduttiddāḷe.}
\end{quote}

\begin{quote}
She is studying archaeology at college.
\end{quote}

\kn{ಪುರಾತತ್ವ} \kn{ಇಲಾಖೆ} is the department that looks after old monuments.

\begin{cousinweb}
\textbf{\kn{ಪುರಾತತ್ವ}} (\emph{purātatva}) is borrowed from Sanskrit, a modern compound of \emph{purā}, \textquotedblleft{}formerly, of old\textquotedblright{} (Monier-Williams), and \emph{tattva}, \textquotedblleft{}the real state, truth\textquotedblright{} (CDIAL 5642).
\end{cousinweb}

\begin{grammarlens}[title={What you've built}]
One more word for telling what happened.
\end{grammarlens}

\subsection*{Your turn}
\begin{itemize}
\raggedright
  \item \emph{Say it:} \emph{purātatva}
  \item \emph{Say it:} \emph{purātatva}, once more
  \item \emph{Say it:} \emph{avaḷu kālējinalli purātatva ōduttiddāḷe.} — She is studying archaeology at college.
  \item \emph{Recall:} say the Kannada for a prince, then the Kannada for an excavation, then say \emph{purātatva} again
\end{itemize}

\begin{tcolorbox}[breakable,colback=teal!4,colframe=teal!35!black,title={Before you move on}]
What does \kn{ಪುರಾತತ್ವ} mean? (\textquotedblleft{}Archaeology\textquotedblright{}.) Say it once more.
\end{tcolorbox}
\section[\texorpdfstring{\kn{ಅವಶೇಷ} (avaśēṣa)}{ಅವಶೇಷ (avaśēṣa)}]{\kn{ಅವಶೇಷ} (avaśēṣa) — ruins; remains}

\label{lesson:KA-C355-avasesa}



\begin{quote}
Before the new one: say the Kannada for an excavation, then the Kannada for archaeology.
\end{quote}

\subsection*{You'll want to know: \kn{ಅವಶೇಷ}}
\textbf{\kn{ಅವಶೇಷ}} — \emph{avaśēṣa} — \textquotedblleft{}ruins; remains\textquotedblright{}.

\begin{quote}
\kn{ಬೆಟ್ಟದ} \kn{ಮೇಲೆ} \kn{ಹಳೆಯ} \kn{ಕೋಟೆಯ} \kn{ಅವಶೇಷಗಳು} \kn{ಇವೆ}.
\end{quote}

\begin{quote}
\emph{beṭṭada mēle haḷeya kōṭeya avaśēṣagaḷu ive.}
\end{quote}

\begin{quote}
There are the ruins of an old fort on the hill.
\end{quote}

Plural: \kn{ಅವಶೇಷಗಳು}, \emph{avaśēṣagaḷu}.

Usually plural: \kn{ಅವಶೇಷಗಳು}, what is left of a building or a city.

\begin{cousinweb}
\textbf{\kn{ಅವಶೇಷ}} (\emph{avaśēṣa}) is borrowed from Sanskrit \emph{avaśeṣa}, \textquotedblleft{}leavings, remainder\textquotedblright{} (Monier-Williams).
\end{cousinweb}

\begin{grammarlens}[title={What you've built}]
One more word for telling what happened.
\end{grammarlens}

\subsection*{Your turn}
\begin{itemize}
\raggedright
  \item \emph{Say it:} \emph{avaśēṣa}
  \item \emph{Say it:} \emph{avaśēṣa}, once more
  \item \emph{Say it:} \emph{beṭṭada mēle haḷeya kōṭeya avaśēṣagaḷu ive.} — There are the ruins of an old fort on the hill.
  \item \emph{Recall:} say the Kannada for an excavation, then the Kannada for archaeology, then say \emph{avaśēṣa} again
\end{itemize}

\begin{tcolorbox}[breakable,colback=teal!4,colframe=teal!35!black,title={Before you move on}]
What does \kn{ಅವಶೇಷ} mean? (\textquotedblleft{}Ruins; remains\textquotedblright{}.) Say it once more.
\end{tcolorbox}
\section[\texorpdfstring{\kn{ಶಾಸನ} (śāsana)}{ಶಾಸನ (śāsana)}]{\kn{ಶಾಸನ} (śāsana) — an inscription; an edict}

\label{lesson:KA-C355-sasana}



\begin{quote}
Before the new one: say the Kannada for archaeology, then the Kannada for ruins.
\end{quote}

\subsection*{You'll want to know: \kn{ಶಾಸನ}}
\textbf{\kn{ಶಾಸನ}} — \emph{śāsana} — \textquotedblleft{}an inscription; an edict\textquotedblright{}.

\begin{quote}
\kn{ಈ} \kn{ಕಲ್ಲಿನ} \kn{ಶಾಸನ} \kn{ಸಾವಿರ} \kn{ವರ್ಷ} \kn{ಹಳೆಯದು}.
\end{quote}

\begin{quote}
\emph{ī kallina śāsana sāvira varṣa haḷeyadu.}
\end{quote}

\begin{quote}
This stone inscription is a thousand years old.
\end{quote}

Plural: \kn{ಶಾಸನಗಳು}, \emph{śāsanagaḷu}.

Old kings cut their orders and gifts into stone: \kn{ಶಿಲಾಶಾಸನ}.

\begin{cousinweb}
\textbf{\kn{ಶಾಸನ}} (\emph{śāsana}) is borrowed from Sanskrit \emph{śāsana}, \textquotedblleft{}a command\textquotedblright{}, and \textquotedblleft{}a royal edict\textquotedblright{} (CDIAL 12420), from \emph{śās}, \textquotedblleft{}to order\textquotedblright{}.
\end{cousinweb}

\begin{grammarlens}[title={What you've built}]
One more word for telling what happened.
\end{grammarlens}

\subsection*{Your turn}
\begin{itemize}
\raggedright
  \item \emph{Say it:} \emph{śāsana}
  \item \emph{Say it:} \emph{śāsana}, once more
  \item \emph{Say it:} \emph{ī kallina śāsana sāvira varṣa haḷeyadu.} — This stone inscription is a thousand years old.
  \item \emph{Recall:} say the Kannada for archaeology, then the Kannada for ruins, then say \emph{śāsana} again
\end{itemize}

\begin{tcolorbox}[breakable,colback=teal!4,colframe=teal!35!black,title={Before you move on}]
What does \kn{ಶಾಸನ} mean? (\textquotedblleft{}An inscription; an edict\textquotedblright{}.) Say it once more.
\end{tcolorbox}
\section[\texorpdfstring{\kn{ಹಸ್ತಪ್ರತಿ} (hastaprati)}{ಹಸ್ತಪ್ರತಿ (hastaprati)}]{\kn{ಹಸ್ತಪ್ರತಿ} (hastaprati) — a manuscript}

\label{lesson:KA-C355-hastaprati}



\begin{quote}
Before the new one: say the Kannada for ruins, then the Kannada for an inscription.
\end{quote}

\subsection*{You'll want to know: \kn{ಹಸ್ತಪ್ರತಿ}}
\textbf{\kn{ಹಸ್ತಪ್ರತಿ}} — \emph{hastaprati} — \textquotedblleft{}a manuscript\textquotedblright{}.

\begin{quote}
\kn{ಗ್ರಂಥಾಲಯದಲ್ಲಿ} \kn{ತಾಳೆಗರಿಯ} \kn{ಹಳೆಯ} \kn{ಹಸ್ತಪ್ರತಿಗಳು} \kn{ಇವೆ}.
\end{quote}

\begin{quote}
\emph{granthālayadalli tāḷegariya haḷeya hastapratigaḷu ive.}
\end{quote}

\begin{quote}
There are old palm-leaf manuscripts in the library.
\end{quote}

Plural: \kn{ಹಸ್ತಪ್ರತಿಗಳು}, \emph{hastapratigaḷu}.

\kn{ಹಸ್ತ}, hand, and \kn{ಪ್ರತಿ}, a copy: a copy made by hand.

\begin{cousinweb}
\textbf{\kn{ಹಸ್ತಪ್ರತಿ}} (\emph{hastaprati}) is borrowed from Sanskrit, a modern compound of \emph{hasta}, \textquotedblleft{}the hand\textquotedblright{} (CDIAL 14024), and \emph{prati}, \textquotedblleft{}towards, against\textquotedblright{} (CDIAL 8540), used in Kannada for \textquotedblleft{}a copy\textquotedblright{}.
\end{cousinweb}

\begin{grammarlens}[title={What you've built}]
One more word for telling what happened.
\end{grammarlens}

\subsection*{Your turn}
\begin{itemize}
\raggedright
  \item \emph{Say it:} \emph{hastaprati}
  \item \emph{Say it:} \emph{hastaprati}, once more
  \item \emph{Say it:} \emph{granthālayadalli tāḷegariya haḷeya hastapratigaḷu ive.} — There are old palm-leaf manuscripts in the library.
  \item \emph{Recall:} say the Kannada for ruins, then the Kannada for an inscription, then say \emph{hastaprati} again
\end{itemize}

\begin{tcolorbox}[breakable,colback=teal!4,colframe=teal!35!black,title={Before you move on}]
What does \kn{ಹಸ್ತಪ್ರತಿ} mean? (\textquotedblleft{}A manuscript\textquotedblright{}.) Say it once more.
\end{tcolorbox}
